%! Characteristics of Quadratic Functions — Unit 3, Lesson 3.0 — Guided Notes (Answer Key)
\documentclass[10pt]{article}
\usepackage{algebra2-article}
\usepackage{algebra2-key}   % pulls in -boxes; do NOT also load -boxes
\newcommand{\TallMath}[1]{$\displaystyle #1\rule[-1.4em]{0pt}{3.2em}$}
\newcommand{\lgrid}[4]{%
  \draw[step=1, linegray!40, very thin] (#1,#3) grid (#2,#4);
  \draw[->, semithick, charcoal] (#1,0)--(#2,0) node[below right, font=\scriptsize]{$x$};
  \draw[->, semithick, charcoal] (0,#3)--(0,#4) node[left, font=\scriptsize]{$y$};}
% Vocab answer for the key: bold forest term, then the filled definition on a write-line.
\newcommand{\vocabans}[2]{%
  \noindent\textbf{\textcolor{forest}{#1:}}\\[1pt]\ansline{#2}}

\begin{document}

\pageheader{Unit 3, Lesson 3.0}{Guided Notes --- Answer Key}
\namedateperiod

\vspace{0.10in}

\begin{objectivebox}
\textbf{By the end of this lesson, I will be able to\dots}
\begin{itemize}\setlength{\itemsep}{2pt}
  \item name a parabola's \emph{vertex} (its \emph{turning point}) and \emph{axis of symmetry}, and say whether the vertex is a \emph{maximum} or a \emph{minimum};
  \item read the \emph{domain}, \emph{range}, \emph{intercepts}, and \emph{zeros} of a quadratic from its graph, and the intervals where it is increasing, decreasing, positive, or negative;
  \item describe a parabola's \emph{end behavior} and use its \emph{symmetry} to justify what I read.
\end{itemize}
\end{objectivebox}

\vspace{0.06in}

\begin{vocabbox}
\small
Fill in each term as we build it together.
\vspace{2pt}
\vocabans{Parabola}{the U-shaped graph of a quadratic function $f(x)=ax^2+bx+c$; opens up if $a>0$, down if $a<0$}
\vocabans{Vertex (turning point)}{the point where the parabola turns; the lowest point (if it opens up) or highest point (if it opens down)}
\vocabans{Axis of symmetry}{the vertical line through the vertex that the parabola is a mirror image across}
\vocabans{Maximum / minimum value}{the $y$-value at the vertex --- the greatest output (max, opens down) or least output (min, opens up)}
\vocabans{End behavior}{what the outputs do as $x\to+\infty$ and $x\to-\infty$; both arms rise if $a>0$, both fall if $a<0$}
\end{vocabbox}

\begin{hookbox}
A ball is thrown upward. Its height after $t$ seconds is $h(t) = -16t^2 + 32t + 48$ feet.
\begin{itemize}\setlength{\itemsep}{1pt}
  \item How high was the ball when it was released (at $t=0$)? \ansline{$48$ feet --- the $y$-intercept $(0,48)$.}
  \item The graph rises, turns, and falls. What is happening to the ball at the \emph{turning point}? \ansline{It is at its highest point --- the vertex $(1,64)$, a maximum height of $64$ ft at $t=1$ s.}
\end{itemize}
The ball's whole flight --- its start, its highest point, and where it lands --- is written in the
\emph{shape} of one curve. That U-shaped (here, upside-down-U) curve is a \textbf{parabola}, and today
we learn to read it.
\end{hookbox}

\begin{notesbox}{1. What Makes a Function Quadratic}
A \textbf{quadratic function} has the form $f(x) = ax^2 + bx + c$ with $a \ne 0$ --- its highest power
is $x^{\textcolor{keyred}{\mathbf{2}}}$. Build a table for $f(x) = x^2 - 2x - 3$:
\begin{center}
\renewcommand{\arraystretch}{1.25}
\begin{tabular}{|c|c|c|c|c|c|c|}
\hline
$x$ & $-2$ & $-1$ & $0$ & $1$ & $2$ & $3$ \\ \hline
$f(x)$ & $5$ & $0$ & $-3$ & \ans{$-4$} & \ans{$-3$} & \ans{$0$} \\ \hline
\end{tabular}
\end{center}
The outputs fall, reach a lowest value, then rise again --- they do \emph{not} change by a constant
amount, so the graph is \textbf{not} a line. Its graph is a symmetric U-shaped curve called a
\ans{parabola}. Because $a = 1 > 0$, this parabola opens \ans{up} (a bowl); if $a<0$ it opens
\ans{down} (a dome).
\end{notesbox}

\begin{notesbox}{2. Reading a Parabola: The Characteristics}
Here is the graph of $f(x) = x^2 - 2x - 3$. Read each feature straight off the picture.
\begin{center}
\begin{tikzpicture}[scale=0.52]
  \lgrid{-3}{5}{-5}{6}
  \draw[navy, dashed, semithick] (1,-5)--(1,6);
  \node[navy, font=\scriptsize, right] at (1.05,5.4){axis $x=1$};
  \begin{scope}
    \clip (-3,-5) rectangle (5,6);
    \draw[very thick, forest, domain=-2.2:4.2, samples=80, smooth]
      plot (\x,{(\x)*(\x) - 2*(\x) - 3});
  \end{scope}
  \fill[charcoal] (1,-4) circle (3pt) node[below right, font=\scriptsize]{vertex $(1,-4)$};
  \fill[charcoal] (0,-3) circle (3pt) node[right, font=\scriptsize]{$(0,-3)$};
  \fill[charcoal] (-1,0) circle (3pt) node[above left, font=\scriptsize]{$(-1,0)$};
  \fill[charcoal] (3,0) circle (3pt) node[above right, font=\scriptsize]{$(3,0)$};
\end{tikzpicture}
\end{center}
\begin{enumerate}[leftmargin=1.6em, itemsep=3pt]
  \item \textbf{Vertex (turning point):} the point where the curve turns, $(\textcolor{keyred}{\mathbf{1}},\,\textcolor{keyred}{\mathbf{-4}})$.
        Here the curve turns from falling to rising, so the vertex is a \textbf{minimum}.
  \item \textbf{Axis of symmetry:} the vertical mirror line through the vertex, $x = \textcolor{keyred}{\mathbf{1}}$. The
        two halves of the parabola are reflections across it.
  \item \textbf{Minimum value:} the \emph{output} at the vertex, $f = \textcolor{keyred}{\mathbf{-4}}$. (A minimum is a
        $y$-value, not a point.)
  \item \textbf{$y$-intercept:} set $x=0$: the point $(0,\,\textcolor{keyred}{\mathbf{-3}})$.
  \item \textbf{$x$-intercepts / zeros:} where $f(x)=0$: $x = \textcolor{keyred}{\mathbf{-1}}$ and $x = \textcolor{keyred}{\mathbf{3}}$. A
        parabola can have \emph{two}, one, or no zeros.
  \item \textbf{Domain / range:} domain is \ans{all reals, $(-\infty,\infty)$}; because the lowest output is $-4$ and the
        curve rises forever, the range is \ans{$y\ge -4$, i.e.\ $[-4,\infty)$}.
  \item \textbf{Increasing / decreasing:} the graph \emph{falls} then \emph{rises}, turning at the
        vertex. Decreasing on \ans{$(-\infty,1)$}; increasing on \ans{$(1,\infty)$}.
  \item \textbf{Positive / negative:} above the $x$-axis (so $f(x)>0$) when $x<-1$ \emph{or} $x>3$;
        below it (so $f(x)<0$) on the interval \ans{$-1<x<3$} (between the zeros).
  \item \textbf{End behavior:} follow both arms. As $x \to +\infty$, $f(x) \to \ans{$+\infty$}$; and as
        $x \to -\infty$, $f(x) \to \ans{$+\infty$}$. Both arms go \emph{up} because $a>0$.
\end{enumerate}
\end{notesbox}

\begin{notesbox}{3. Opening Direction \& Symmetry}
The sign of $a$ decides which way a parabola opens --- and therefore whether the vertex is a
\textbf{minimum} or a \textbf{maximum}.
\begin{center}
\begin{tikzpicture}[scale=0.44]
  % Parent y = x^2 : even symmetry about the y-axis
  \begin{scope}
    \lgrid{-3}{3}{-1}{6}
    \draw[navy, dashed, semithick] (0,-1)--(0,6);
    \begin{scope}
      \clip (-3,-1) rectangle (3,6);
      \draw[very thick, forest, domain=-2.4:2.4, samples=60, smooth] plot (\x,{(\x)*(\x)});
    \end{scope}
    \fill[charcoal] (0,0) circle (3pt);
    \node[charcoal, font=\scriptsize] at (0,-2.2){$y=x^2$ \; (opens up, min)};
  \end{scope}
  % g(x) = -x^2 + 2x + 3 : opens down, max
  \begin{scope}[xshift=10cm]
    \lgrid{-3}{5}{-4}{6}
    \draw[navy, dashed, semithick] (1,-4)--(1,6);
    \begin{scope}
      \clip (-3,-4) rectangle (5,6);
      \draw[very thick, forest, domain=-2.2:4.2, samples=80, smooth]
        plot (\x,{-1*(\x)*(\x) + 2*(\x) + 3});
    \end{scope}
    \fill[charcoal] (1,4) circle (3pt);
    \node[charcoal, font=\scriptsize] at (1,-5.2){$g(x)=-x^2+2x+3$ \; (opens down, max)};
  \end{scope}
\end{tikzpicture}
\end{center}
\begin{itemize}[leftmargin=1.5em, itemsep=2pt]
  \item \textbf{Even symmetry.} The parent $y=x^2$ is a mirror image across the \ans{$y$}-axis:
        inputs $x$ and $-x$ give the \emph{same} output, e.g.\ $(-2)^2 = (2)^2 = \textcolor{keyred}{\mathbf{4}}$.
  \item \textbf{Every parabola} is symmetric across its \textbf{axis of symmetry}. For $g$ above, the
        axis is $x=\textcolor{keyred}{\mathbf{1}}$, so points the same distance left and right of it share an output.
  \item \textbf{$g$ opens down} ($a=-1<0$), so its vertex $(1,4)$ is a \textbf{maximum}; its maximum
        value is \ans{$4$}, its range is \ans{$y\le 4$}, and \emph{both} end-behavior arms go
        \ans{down ($-\infty$)}.
\end{itemize}
\end{notesbox}

\begin{practicebox}
The graph below is $p(x) = x^2 - 4$. Read each characteristic and write it on the line.
\begin{center}
\begin{tikzpicture}[scale=0.5]
  \lgrid{-4}{4}{-5}{5}
  \draw[navy, dashed, semithick] (0,-5)--(0,5);
  \begin{scope}
    \clip (-4,-5) rectangle (4,5);
    \draw[very thick, forest, domain=-3:3, samples=70, smooth] plot (\x,{(\x)*(\x) - 4});
  \end{scope}
  \fill[charcoal] (0,-4) circle (3pt) node[below right, font=\scriptsize]{$(0,-4)$};
  \fill[charcoal] (-2,0) circle (3pt) node[above left, font=\scriptsize]{$(-2,0)$};
  \fill[charcoal] (2,0) circle (3pt) node[above right, font=\scriptsize]{$(2,0)$};
\end{tikzpicture}
\end{center}
\begin{enumerate}[leftmargin=1.6em, itemsep=3pt]
  \item Vertex: \ans{$(0,-4)$} \quad Min or max? \ans{minimum} \quad Axis of symmetry: $x=$ \ans{$0$}
  \item Zeros: \ans{$x=-2$ and $x=2$} \qquad $y$-intercept: \ans{$(0,-4)$}
  \item Domain: \ans{all reals} \qquad Range: \ans{$y\ge -4$}
  \item Decreasing on \ans{$(-\infty,0)$}; \; increasing on \ans{$(0,\infty)$}; \; negative on \ans{$-2<x<2$}.
\end{enumerate}
\end{practicebox}

\begin{teachernote}
The whole lesson is one move --- \emph{read a feature off the picture} --- now on a curve that turns.
Three predictable errors: (1) giving the min/max as a point instead of the $y$-value (item~3 in \S2);
(2) fusing ``decreasing'' with ``negative'' (the parabola is decreasing on $(-\infty,1)$ but only
negative on $(-1,3)$); and (3) writing the vertex $x$-coordinate as the range. \S3 exists to lock in
that the \emph{sign of $a$} controls opening direction, min-vs-max, and end behavior all at once. Keep
returning to the axis of symmetry as the justification tool: equal steps left and right share an
output.
\end{teachernote}

\end{document}
